\phantomsection
\addcontentsline{toc}{chapter}{Editor's Preface}
\fancyhead[C]{{\LARGE Editor's Preface}}
\par\noindent
\begin{quoting}
	The Lord hath hallowed his tabernacle: for this is the house of the Lord, wherein men may call upon his Name; whereof it is written, My Name shall be therein, saith the Lord.
\end{quoting}

\lett{W}{ith} great joy and evangelical jubilation, Apologia Anglicana is blest to release this edition of the Book of the Common Prayer in the Year of Our Lord 2027. On this day, the \nth{477} anniversary of the first Book of Common Prayer in 1549, we remember the glorious liturgical heritage of our ancestors. And in our time, we celebrate our redemption from heresy and schism into the Orthodox Church.\\

What the English Reformers saw dimly, we rejoice in the fulness of truth. We pray for the descendants of those blessed isles, that they find their home and rest within the mystical Body of Christ. We are intensely thankful for all Orthodox Christians and especially our God-fearing bishops who gave not up on us nor harshly imposed a foreign rite upon us, which our sins would most certainly have deserved. But rather, with patient care for our souls, we rejoice in thanksgiving for not only retaining our liturgical heritage but having it healed from all error.\\

So now, moving forward, we worship the Lord in the beauty of holiness. We offer right worship, knowing him whom we worship. And we beseech the Lord every day, in every hour, for his grace to bring us unto himself, the only wise God, world without end. Amen.\\

\rightline{On the Feast of Whitsunday}

\vspace{1\baselineskip}

\rightline{\emph{Pax Christi,}}

\vspace{1\baselineskip}

\rightline{Mr. Augustine Watson, MSt}

\rightline{\emph{Chief Editor, Apologia Anglicana}}